\documentclass[withoutpreface]{cumcmthesis}
\usepackage{url} % 支持网址链接，便于引用网络资源
\usepackage{graphicx} % 插入和管理图片
\usepackage{pdflscape} % 使页面横向显示，适合宽表格
\usepackage{tabularx} % 支持自动调整宽度的表格
\usepackage{afterpage} % 控制内容在下一页显示
\usepackage{longtable} % 支持跨页长表格
\usepackage{booktabs}
\usepackage[numbers,sort&compress]{natbib} % 参考文献排序与压缩显示

\usepackage[figuresright]{rotating} % 支持图表旋转


\usepackage{amsmath} % 数学公式增强
\usepackage{amssymb} % 更多数学符号
\usepackage{geometry} % 页面边距自定义
\usepackage{enumitem} % 列表样式自定义
\usepackage{caption} % 图表标题样式自定义
\usepackage{fancyhdr} % 页眉页脚自定义
\usepackage{titlesec} % 章节标题样式自定义

\usepackage{listings} % 插入和高亮显示代码
\newfontfamily\codefont{DejaVu Sans Mono}
% 设置代码样式
\lstset{
	numbers             =   left,               % 行号显示在左侧
	showspaces          =   false,              % 不显示空格符号
	numberstyle         =   \zihao{-5}\codefont,% 行号风格，小五号，等宽字体
	showstringspaces    =   false,              % 不显示字符串中的空格
	captionpos          =   t,                  % 标题位置在顶部
	frame               =   lrtb,               % 显示四周边框
	breaklines          =   true,               % 自动换行
	columns             =   fixed,              % 固定列宽
	basewidth           =   0.5em,              % 基础宽度
	caption             =   \raggedright,       % 标题左对齐
	literate            =   {ρ}{{$\rho$}}1
	                        {σ}{{$\sigma$}}1
	                        {δ}{{$\delta$}}1
	                        {κ}{{$\kappa$}}1,    % 保证附录代码中的希腊字母可正确显示
}

% Python 代码样式
\lstdefinestyle{Python}{
	language        =   Python, % 语言设为Python
	basicstyle      =   \zihao{-5}\codefont\color{black}, % 基本风格，小五号，等宽字体，黑色
	keywordstyle    =   \color{blue}\bfseries\itshape, % 关键字风格，蓝色加粗斜体
	stringstyle     =   \color{magenta}, % 字符串风格，洋红色
	commentstyle    =   \color{red}\itshape, % 注释风格，红色斜体
}

% MATLAB 代码样式
\lstdefinestyle{MATLAB}{
	language        =   Matlab, % 语言设为MATLAB
	basicstyle      =   \zihao{-5}\codefont\color{black}, % 基本风格，小五号，等宽字体，黑色
	keywordstyle    =   \color{green}\bfseries\itshape, % 关键字风格，绿色加粗斜体
	stringstyle     =   \color{purple}, % 字符串风格，紫色
	commentstyle    =   \color{gray}\itshape, % 注释风格，绿色斜体
}



\usepackage{xcolor} % 颜色支持，代码和表格可用
\usepackage{multicol} % 多栏排版

\usepackage{booktabs} % 优美的表格线
\usepackage{hyperref} % 支持PDF跳转、超链接
\usepackage{subcaption} % 子图标题，多个图片并排显示
\usepackage {float}  % 更灵活地控制浮动体（如表格、图片）位置
\usepackage{array} % 增强表格功能
\hypersetup{hidelinks,
	colorlinks=true,
	allcolors=black,
	pdfstartview=Fit,
	breaklinks=true
}
\usepackage{siunitx} % 科学计数法和单位格式化
\usepackage{threeparttable}
\usepackage{placeins}
\usepackage{microtype}
\geometry{left=2.4cm,right=2.4cm,top=2.35cm,bottom=2.35cm}
\captionsetup{font=small,labelfont=bf}

\title{锂离子电池快充寿命的建模预测与策略优化研究}
\tihao{A}
\baominghao{}
%\schoolname{BIT}			%学校  
%\yearinput{2025}			%年份  
%\monthinput{09}			%月份
%\dayinput{03}				%日期

\begin{document}

\pagenumbering{arabic}

\maketitle
\begin{center}
    {\heiti\zihao{-4} 摘\quad 要}
\end{center}

锂离子电池快速充电需要兼顾充电效率、寿命衰减与运行安全等问题，但早期循环数据普遍存在寿命终点不可见、策略参数耦合和重复样本非独立等问题，容易造成寿命误判与策略过度外推。
本文基于 MIT--Stanford 电池循环老化数据集，构建“\textbf{右删失识别--稳健退化表征--Hybrid 寿命预测--风险约束优化}”的建模框架，分析快充参数对电池老化的影响并优化充电策略。

\textbf{针对问题一},在真实寿命均不可观测的条件下构建早期老化代理指标。由于所有电池 SOH 均高于 \textbf{80\% 失效阈值}，循环观测次数不能等效为真实寿命。以第 20--30 次循环的中位 SOH 为稳定基准，经后续数据验证选取第 80--150 次循环的 \textbf{Theil--Sen 稳健退化率}表征早期衰减，并利用第 141--150 次循环的端点相对 SOH 损失复核策略排序。结果表明，S9、S8 分别为同一实验层内典型低、高衰减策略；S8 相比 S9 的退化率增加 \textbf{1.773 个百分点/100 次}，而充电时间仅缩短 \textbf{0.320--0.942 min}。

\textbf{针对问题二},以六个独立策略点为正式推断单位，采用 Kruskal--Wallis 检验、精确 Spearman 相关、Holm 校正及整策略留出 Ridge 模型分析参数作用。六策略退化率存在总体差异（\textbf{\(p=0.0130\)}）；仅删除 S8 后 \(p\) 升至 \textbf{0.1173}，说明差异主要由 S8 驱动。阶段切换 SOC 参数 \(Q_1\) 与策略中位退化率的相关系数为 \textbf{\(\rho=-0.886\)}，原始 \(p=\textbf{0.0333}\)，但三参数校正后 \(p=\textbf{0.100}\)。因此，\(Q_1\) 仅可视为具有稳定方向性的重点关注参数，现有数据尚不能确认三个快充参数的独立效应及严格因果关系。

\textbf{针对问题三},分别完成第 151--200 次循环的短期 SOH 预测和 80\% SOH 寿命终点的远期外推。构建融合 \textbf{PCA--Ridge 数据驱动模型与拉伸指数经验模型的 Hybrid 混合预测模型}，并采用 Outer/Inner 双层嵌套交叉验证完成模型选择与泛化评估。Hybrid 模型平均 \textbf{RMSE 为 0.0480 个百分点}，较内层最优单模型降低 \textbf{5.7\%}，点覆盖率和整轨迹覆盖率分别达到 \textbf{98.8\% 和 92.5\%}。对于远期寿命，则采用三类单调退化模型构建跨模型外推包络，降低单点外推造成的寿命误判风险。

\textbf{针对问题四},以实测充电时间和早期退化率为双目标，建立由双分支max-Guard、Outer LOPO误差下限和凸包—最近邻可信域共同约束的风险Pareto模型。结合报告精度下的$\varepsilon$-Pareto判定与20\,000次电池级Bootstrap，推荐\textbf{S5（5.3C--54\%--4.0C）作为速度与寿命的均衡代表}，推荐\textbf{S9（4.8C--80\%--4.8C）作为健康优先策略}。连续Balanced方案（5.480C--22.39\%--4.656C）相对S5仅预测快$0.108\,\mathrm{min}$，退化率却增加35.1\%；相对S9预测快约$1.103\,\mathrm{min}$，退化率增加约81.1\%。因此，本文将\textbf{Balanced（5.480C--22.39\%--4.656C）作为兼顾充电时间与寿命衰减的连续优化备选策略}，并建议经独立退化实验验证后再行应用。

\noindent\textbf{关键词：}锂离子电池；快充策略；右删失；Hybrid 寿命预测；Pareto 优化


\clearpage

\section{问题重述}
\subsection{问题背景}
锂离子电池快充需要兼顾补能时间与容量衰减，而电流引起的极化、温升和副反应又随SOC区间变化。附件源自MIT--Stanford公开老化实验，相关研究表明早期循环特征可预测后续性能，并可服务于闭环快充优化\cite{severson2019,attia2020}。本题的关键是在有限、非正交的策略设计中区分可靠证据与外推结论。

\subsection{问题重述}
\textbf{问题1}要求整理电池信息与寿命数据，比较策略分布，提取典型长、短寿命策略并解释差异。首先必须判断80\% SOH终点是否可见，再为不可见终点建立可验证的早期代理。

\textbf{问题2}要求检验策略差异，分析$C_1,Q_1,C_2$及SOC区间与衰减的关系，并讨论影响程度和可能机制。三参数并非独立设计，需区分重复电池与独立策略点。

\textbf{问题3}要求用前150次数据预测9块测试电池第151--200次SOH，外推80\% EOL，并评价早期数据长度与策略特征的作用。

\textbf{问题4}要求同时降低充电时间与寿命衰减，在有实验支撑的参数域内给出推荐并量化外推风险。




\section{问题分析}
非测试电池记录到200次，9块测试电池只记录到150次；最小平滑SOH为94.97\%，故全部80\%寿命事件右删失。由此确定三个分析原则：
\begin{enumerate}[label=(\arabic*)]
  \item 问题一在共同150次窗口内构造早期退化代理，不把观测截止次数当作真实寿命；
  \item 问题三以40块非测试电池的第151--200次记录回测，全部选择步骤限于Outer训练端；
  \item EOL与连续策略优化均属于外推，必须报告不确定性并限制适用范围。
\end{enumerate}

问题一依次完成信息整理、分布统计、典型策略提取和差异解释。\texttt{dataset\_id}为原始分组字段，本文仅将其作为批次/区组代理；“新结构”直接取自\texttt{policy}字段的\texttt{NEWSTRUCTURE}后缀。为降低可见混杂，典型策略优先在\texttt{dataset\_id}=3且结构标签一致的子集中比较。

问题二以$r_i$为主响应、$D_i$为替代响应，正式推断基于六个独立策略点，重复电池仅估计策略内变异。问题三将物理参数与实验元数据分开，并解耦短期预测和EOL外推。问题四把S4--S9聚合为策略级响应，以Outer LOPO评估Guard，并在凸包—最近邻可信域内搜索Pareto候选；泛化不足时只推荐已有实测策略。


\section{模型假设}
\begin{enumerate}[label=\textbf{A\arabic*.}]
  \item 附件中的\texttt{dataset\_id}作为原始实验分组字段，可用于构造内部比较区组；因题目说明未给出其完整物理语义，本文将其视为批次/区组代理，而不把组间差异解释为已确认的结构因果效应。结构标签仅按原始\texttt{policy}是否含\texttt{NEWSTRUCTURE}识别。
  \item 平滑SOH是容量健康状态的主要观测量；短时容量回弹属于测量与弛豫波动，不代表永久性“负衰减”。
  \item 前150次内总体退化方向近似单调，允许局部噪声；未来50次预测沿用该早期机理，而80\% EOL不保证沿用同一速率。
  \item 问题四的连续优化仅适用于S4--S9六个策略点的参数凸包及最近邻可信域；域外不作推荐，Guard预测不替代实测证据。
\end{enumerate}
\section{主要符号说明}
\begin{table}[H]
\centering
\caption{主要符号及含义}
\label{tab:main_symbols}
\small
\setlength{\tabcolsep}{2pt}
{ % 分组，只作用于本tabularx
\renewcommand{\arraystretch}{0.95} % ← 调这个数值缩小行高，越小越紧凑
\begin{tabularx}{\textwidth}{cXcX}
\toprule
符号 & 含义 & 符号 & 含义 \\
\midrule

$S^c_{i,n}$
& 电池$i$在第$n$次循环的清洗后SOH
& $B_i$
& 电池$i$在Cycle 20--30内的稳定SOH基准 \\

$R_{i,n}=S^c_{i,n}/B_i$
& 相对于稳定基准归一化后的SOH
& $r_i,\ D_i$
& Cycle 80--150稳健退化率与Cycle 141--150端点相对损失 \\

$L_i,\ C_i,\ \widetilde L_i,\ \delta_i$
& 潜在80\% SOH寿命、删失时间、可观测时间及寿命事件指示量
& $(C_{1,p},Q_{1,p},C_{2,p})$
& 策略$p$的第一阶段倍率、阶段切换SOC和第二阶段倍率 \\

$E_{p,k},\ E_{L,p},\ E_{H,p}$
& 策略$p$的分区间、低SOC及高SOC倍率暴露
& $\tau_i^{(150)},\ T_p$
& 电池前150次循环的充电时间中位数及策略级实测充电时间 \\

$\ell,\ \boldsymbol{x}_i^{(\ell)}$
& 早期观测长度及由前$\ell$次循环构造的多源特征向量
& $\widehat S_{i,n}$
& 电池$i$在第$n$次循环的SOH预测值 \\

$N_{i,\mathrm{EOL}}^{(m)},\
 \widetilde N_{i,\mathrm{EOL}},\
 R_{\mathrm{ext},i}$
& 模型族$m$下的EOL样本、跨模型中位中心及连续外推指数
& $r_{p,\mathrm{rep}},\
 Y_p=\ln(r_{p,\mathrm{rep}}/10^{-5})$
& 策略$p$的代表退化率及其对数健康响应 \\

$T_{\mathrm{CC}},\ \mu_T,\ \sigma_T$
& 理论充电时间及时间Guard的预测均值和不确定性
& $\mu_Y,\ \sigma_Y$
& 健康Guard的对数退化率预测均值和不确定性 \\

$\mathcal X_{\mathrm{safe}},\
 d_{\min},\ \delta_0$
& 连续策略可信域、标准化最近邻距离及其阈值
& $F_T,\ F_Y,\ \mathcal P$
& 风险修正的时间目标、健康目标及Pareto非支配集 \\

\bottomrule
\end{tabularx}
}%分组结束，arraystretch自动恢复原来的值
\end{table}
\section{模型的建立与求解}

\subsection{问题一：早期寿命表征与典型策略比较}

\subsubsection{数据整理与寿命可观测性}
以$(\text{battery\_id},\text{cycle})$为主键，共得到49块电池、9类策略和9350条循环记录。原始列完整保留，在clean副本中仅修正电池1第12次容量/SOH孤立尖峰以及电池2、3第12次零内阻；S2的$C_1$空值源于0--80\% SOC单段3.6C协议，不作一般缺失值插补。平滑轨迹只用于展示，全部指标均由清洗序列$S^c_{i,n}$计算。

形成/激活阶段会造成初始SOH偏移，故以Cycle 20--30中位数为稳定基准：
\begin{equation}
B_i=\operatorname{median}_{20\le n\le30}S^c_{i,n},\qquad
R_{i,n}=\frac{S^c_{i,n}}{B_i}.
\label{eq:relative-soh}
\end{equation}
原始SOH及Cycle 15--25、20--30、25--35三种基准窗所得策略排序的Spearman相关均为1.000，说明归一化不改变主要排序。图\ref{fig:p1strategy}展示策略中位轨迹及电池间四分位区间。

\begin{figure}[htbp]
\centering
\includegraphics[width=0.50\textwidth]{texfile/figures/fig_p1_strategy_trajectories.pdf}
\caption{九类策略的稳定基准归一化SOH轨迹}
\label{fig:p1strategy}
\end{figure}

令潜在80\% SOH寿命为$L_i=\inf\{n:S_{i,n}\le0.8\}$，删失时间$C_i$对测试、非测试电池分别为150和200，则
\begin{equation}
\widetilde L_i=\min(L_i,C_i),\qquad \delta_i=\mathbb I(L_i\le C_i).
\label{eq:observed-life}
\end{equation}
全部电池清洗SOH$S^c_{i,n}$最小值约94.66\%，绘图平滑序列$S^{sm}_{i,n}$最小值94.97\%，二者均远高于80\%失效阈值，全部寿命事件满足右删失条件$\delta_i=0$,只能得到$L_i>C_i$。逐电池完整索引见支撑材料；表\ref{tab:p1inventory}给出策略级汇总。

\begin{table}[htbp]
\centering
\caption{电池、策略、寿命下界与充电时间索引}
\label{tab:p1inventory}
\footnotesize
\begin{tabularx}{\textwidth}{ccXcc}
\toprule
策略 & 电池数 & 快充形式 & 80\% SOH寿命信息 & 充电时间中位数/min\\
\midrule
S1&3&3.6C--80\%--3.6C&1块$>150$，2块$>200$&13.342\\
S2&3&单段3.6C（0--80\% SOC）&1块$>150$，2块$>200$&13.342\\
S3&3&4.8C--80\%--4.8C&1块$>150$，2块$>200$&10.054\\
S4&7&5.0C--67\%--4.0C&1块$>150$，6块$>200$&10.043\\
S5&8&5.3C--54\%--4.0C&1块$>150$，7块$>200$&10.043\\
S6&8&5.6C--19\%--4.6C&1块$>150$，7块$>200$&10.043\\
S7&8&5.6C--36\%--4.3C&1块$>150$，7块$>200$&10.043\\
S8&3&3.7C--31\%--5.9C&1块$>150$，2块$>200$&10.096\\
S9&6&4.8C--80\%--4.8C&1块$>150$，5块$>200$&11.038\\
\bottomrule
\end{tabularx}
\end{table}

\subsubsection{早期退化表征与策略分布}
由于没有失效事件，各策略真实寿命中位数均不可辨识。为比较早期寿命表现，使用40块具有第151--200次观测的非测试电池验证退化窗口。候选窗口同时满足Pearson相关不低于0.96、Spearman相关不低于0.80时，选择覆盖最长者；分别比较各候选窗口退化斜率与Cycle 150–200后续真实退化斜率的Pearson相关系数、Spearman秩相关系数及斜率MAE，结果如图3所示。Cycle 80--150、90--150和100--150均满足条件，故以含71个循环点的80--150为主窗口，数值最优的100--150用于敏感性检验。起点在60--100间变化时，S9/S1持续位于低退化端，S8持续最高。

\begin{figure}[htbp]
\centering
\includegraphics[width=0.50\textwidth]{texfile/figures/fig_p1_window_validation.pdf}
\caption{候选退化窗口的后续观测验证}
\label{fig:p1window}
\end{figure}

主退化率与端点损失分别定义为
\begin{equation}
r_i=-\operatorname{median}_{80\le u<v\le150}
\frac{R_{i,v}-R_{i,u}}{v-u},
\label{eq:theilsen-rate}
\end{equation}
\begin{equation}
D_i=100\left[1-\operatorname{median}_{141\le n\le150}R_{i,n}\right].
\label{eq:endpoint-loss}
\end{equation}
$r_i$刻画后期下降速度，$D_i$刻画公共终点累计损失；两者源于同一序列，不视为统计独立验证。\textbf{$r_i$和$D_i$越大，均表示早期寿命表现越差}。九类快充策略下两项寿命代理指标的电池级分布分别见
图\ref{fig:p1rate}和图\ref{fig:p1endpoint}，
其中散点表示独立电池，箱体表示第一至第三四分位区间，
箱内竖线表示中位数。

\begin{figure}[htbp] % 改动这里，去掉!b
\centering
\begin{minipage}[b]{0.47\textwidth}
\centering
\includegraphics[width=\textwidth]{texfile/figures/fig_p1_rate_distribution.pdf}
\captionof{figure}{各策略Cycle 80--150相对SOH稳健退化率$r_i$的分布}
\label{fig:p1rate}
\end{minipage}
\hfill
\begin{minipage}[b]{0.47\textwidth}
\centering
\includegraphics[width=\textwidth]{texfile/figures/fig_p1_endpoint_distribution.pdf}
\captionof{figure}{各策略Cycle 141--150端点相对SOH损失$D_i$的分布}
\label{fig:p1endpoint}
\end{minipage}
\end{figure}

共同前150次循环的电池级充电时间及策略汇总定义为
\begin{equation}
\tau_i^{(150)}=\operatorname{median}_{1\le n\le150}t_{i,n},\qquad
\tau_p^{\rm cycle}=\operatorname{median}_{i\in p}\tau_i^{(150)},
\label{eq:charge-time-cycle}
\end{equation}
并以汇总表平均充电时间的策略中位数$\tau_p^{\rm summary}$作敏感性口径。表\ref{tab:p1strategy}按主退化率排序。

\begin{table}[H]
    \centering
    \caption{九类策略的早期寿命代理与充电时间}
    \label{tab:p1strategy}
    \footnotesize
    \begin{threeparttable}
    \setlength{\tabcolsep}{5pt}
    \renewcommand{\arraystretch}{0.40}
    \begin{tabular}{ccccccc}
        \toprule
        策略 & $n$ & 快充形式 & $10^5 r_i$/cycle & $D_i$/\% & $\tau_{\rm cycle}$/min & $\tau_{\rm summary}$/min \\
        \midrule
        S9 & 6 & 4.8C--80\%--4.8C & 1.997  & 0.182 & 11.038 & 11.158 \\
        S1 & 3 & 3.6C--80\%--3.6C & 2.131  & 0.199 & 13.342 & 13.370 \\
        S5 & 8 & 5.3C--54\%--4.0C & 2.676  & 0.231 & 10.043 & 10.154 \\
        S4 & 7 & 5.0C--67\%--4.0C & 2.822  & 0.238 & 10.043 & 10.165 \\
        S7 & 8 & 5.6C--36\%--4.3C & 2.990  & 0.289 & 10.043 & 10.212 \\
        S6 & 8 & 5.6C--19\%--4.6C & 3.081  & 0.343 & 10.043 & 10.242 \\
        S2 & 3 & 单段3.6C & 3.321  & 0.325 & 13.342 & 13.731 \\
        S3 & 3 & 4.8C--80\%--4.8C & 11.265 & 0.965 & 10.054 & 13.145 \\
        S8 & 3 & 3.7C--31\%--5.9C & 19.732 & 2.329 & 10.096 & 10.838 \\
        \bottomrule
    \end{tabular}
    \def\TPTnoteSettings{%
        \setlength\leftmargin{0pt}%
        \setlength\labelwidth{0pt}%
        \setlength\labelsep{0pt}%
        \rightskip 0pt \leftskip 0pt
    }
    \begin{tablenotes}
        \footnotesize
        \item[] 注：充电时间均先取单电池循环中位数，再聚合为策略中位数，复用同一数据源。
    \end{tablenotes}
    \end{threeparttable}
\end{table}

\subsubsection{典型低、高退化策略识别}
典型策略须满足主、替代指标方向一致，并优先在同一原始分组和结构标签内比较，且每类不少于3块电池。在\texttt{dataset\_id}=3的六类新结构策略中，S9的$r_i$和$D_i$均最低，S8两项均最高，故分别作为典型长、短寿命策略。为识别电池 SOH 的退化规律并确定后续分析区间，本文比较了 S8 与 S9 电池相对健康基准 SOH 随循环次数的演化过程。如图5所示，S9 的 SOH 整体较为稳定，而 S8 在约 20 次循环后呈持续下降趋势。S1虽处于低退化端，但缺少同实验层对照，仅保留为参考。

\begin{figure}[htbp]
\centering
\includegraphics[width=0.70\textwidth]{texfile/figures/fig_p1_typical_comparison.pdf}
\caption{S9与S8的相对SOH轨迹}
\label{fig:p1typical}
\end{figure}
\newpage
对两类电池进行非参数Bootstrap，结果见表\ref{tab:p1bootstrap}。S8样本量仅为3，区间只描述现有重复电池波动，不能替代新增实验。
\begin{table}[htbp]
\centering
\caption{S9与S8早期衰减指标的Bootstrap区间}
\label{tab:p1bootstrap}
\small
\begin{tabular}{ccccc}
\toprule
策略 & $n$ & $r$（百分点/100次） & $D$/\% & $\tau_{\rm cycle}$/min\\
\midrule
S9&6&0.200 [0.137, 0.301]&0.182 [0.129, 0.278]&11.038\\
S8&3&1.973 [1.510, 2.256]&2.329 [1.852, 2.552]&10.096\\
\bottomrule
\end{tabular}
\end{table}

两样本的Cliff效应量定义为
\begin{equation}
\delta_{\rm Cliff}=\frac{1}{mn}\sum_{i=1}^{m}\sum_{j=1}^{n}
\left[\mathbb I(X_i>Y_j)-\mathbb I(X_i<Y_j)\right].
\label{eq:cliff}
\end{equation}
取$X$为S8、$Y$为S9的退化率，得到$\delta_{\rm Cliff}=1.00$，即现有样本完全分离。

\subsubsection{典型策略差异与机理解释}
S8相较S9的退化率增加1.773个百分点/100次，端点损失增加2.147个百分点；共同150次口径下充电时间缩短0.942 min，汇总表口径下缩短0.320 min，Bootstrap区间分别为$[-0.947,-0.927]$和$[-0.369,-0.242]$ min。两种口径均表明S8更快，但收益幅度依赖时间定义。S8、S9差异如表\ref{tab:p1typicaldiff}所示。

\begin{table}[htbp]
\centering
\caption{典型低、高早期衰减策略比较}
\label{tab:p1typicaldiff}
\small
\begin{tabularx}{\textwidth}{p{2.7cm}XXp{3.2cm}}
\toprule
维度 & S9 & S8 & 差异\\
\midrule
协议&4.8C--80\%--4.8C&3.7C--31\%--5.9C&S8在31--80\% SOC高1.1C\\
时间/min&11.038&10.096&S8快0.942 min（8.53\%）\\
温度中位数/$^\circ$C&33.124&35.079&S8高1.955$^\circ$C\\
退化率/（百分点/100次）&0.200&1.973&S8增加1.773\\
端点损失/\%&0.182&2.329&S8增加2.147\\
\bottomrule
\end{tabularx}
\end{table}

S8在31--80\% SOC采用5.9C，较S9高1.1C，并伴随更高温度。较大电流可能增强焦耳热、极化及中高SOC区间的析锂和界面副反应风险\cite{koehler2019,operando2022}。由于题目未观测负极电位、析锂或SEI变化，且两策略同时改变三个参数，该机理仅解释为与数据一致的条件关联。综上，S8以不足1 min的时间收益伴随明显更高的早期衰减；“长寿命”和“短寿命”均指早期相对表现，而非已观测的80\% SOH寿命。
\FloatBarrier


\subsection{问题二：快充参数与早期衰减的分层证据分析}
问题二沿用Cycle 80--150退化率$r_i$，以端点损失$D_i$作替代响应，并用$Y_i=\ln(r_i/10^{-5})$稳定联合回归的残差尺度。同一策略内电池共享$C_1,Q_1,C_2$，重复电池只估计策略内变异，为避免不同工况及结构标签造成的混杂影响，本文以 \texttt{dataset\_id} 为基本分组单位，并结合电池编号与 \texttt{NEWSTRUCTURE} 标签确定各组数据的分析用途。全9种策略用于描述性比较；正式参数推断限制在\texttt{dataset\_id}=3的S4--S9六个策略点。该字段仅作为原始批次/区组代理，不解释为已确认的结构原因。

\begin{table}[htbp]
\centering
\caption{原始分组字段及分析用途}
\label{tab:dataset-source}
\small
\renewcommand{\arraystretch}{1.25}
\setlength{\tabcolsep}{6pt}

\begin{tabularx}{\textwidth}{ccp{5.2cm}X}
\toprule
\texttt{dataset\_id} & 电池数 & 策略与标签 & 分析用途\\
\midrule

1 & 3 & S1，无\texttt{NEWSTRUCTURE}
& 低倍率参考，不作组内参数推断\\

2 & 6 & S2、S3，无\texttt{NEWSTRUCTURE}
& 展示跨组差异\\

3 & 40 & S4--S9，均含\texttt{NEWSTRUCTURE}
& 六策略点精确推断\\

\bottomrule
\end{tabularx}
\end{table}



\subsubsection{不同快充策略的总体差异}
各策略样本量仅3--8且分布偏态，采用Kruskal--Wallis秩检验。表\ref{tab:p2kw}给出了不同样本范围下的检验统计量 (H)、显著性 (p) 值及效应量 ($\varepsilon_H^2$),且如图\ref{fig:p2dist}所示，S8 策略的主退化率显著高于其余策略，表明其退化程度更为严重。。\texttt{dataset\_id}=3内主退化率$r_i$的$p=0.0130$，以$D_i$复核为$p=0.0111$。由此可知，\texttt{dataset\_id}=3内六种策略的寿命差异是存在显著总体差异。逐策略删除时，发现仅删除S8会使总体差异消失（$p=0.1173$），
删除其余策略所得$p$为0.0080--0.0297，说明现有数据支持“存在一个明显不利策略”，而非其余五种策略普遍可分。

\begin{table}[htbp]
\centering
\caption{策略总体差异检验}
\label{tab:p2kw}
\small
\begin{tabular}{lccccc}
\toprule
范围 & 响应 & $K$ & $N$ & $H$ & $p$；$\epsilon_H^2$\\
\midrule
全9策略 & $r_i$ & 9 & 49 & 24.246 & 0.00208；0.406\\
\texttt{dataset\_id}=3 & $r_i$ & 6 & 40 & 14.445 & 0.01302；0.278\\
\texttt{dataset\_id}=3（删S8） & $r_i$ & 5 & 37 & 7.375 & 0.11733；0.105\\
\texttt{dataset\_id}=3 & $D_i$ & 6 & 40 & 14.844 & 0.01105；0.290\\
\bottomrule
\end{tabular}
\end{table}

\begin{figure}[htbp]
\centering
\includegraphics[width=0.67\textwidth]{texfile/figures/fig_p2_dataset3_distribution.pdf}
\caption{S4--S9的电池级主退化率分布}
\label{fig:p2dist}
\end{figure}
\newpage
S8--S9的精确Mann--Whitney原始$p=0.0238$，15组两两比较Holm校正后$p=0.238$；$\delta_{\rm Cliff}=1.00$表明现有样本完全分离，因此可以确认六策略总体存在差异并识别S8为明显不利策略，但不能声称S8与某一具体策略已经通过多重校正后的正式两两显著性检验。
渐近Dunn--Holm只作为方法敏感性，完整矩阵见支撑材料。

\subsubsection{充电参数的边际关联}
以六个策略中位退化率为响应，枚举全部$6!=720$种排列，对统计量$T$定义双侧精确概率
\begin{equation}
p_{\rm exact}=\frac{\#\{\pi:|T_\pi|\ge |T_{\rm obs}|\}}{720}.
\label{eq:p2-exact-permutation}
\end{equation}
由公式（7）计算可得，结果详情见表\ref{tab:spearman_relation}。
\begin{table}[htbp]
    \centering
    \caption{参数与早期退化率的 Spearman 相关关系}
    \label{tab:spearman_relation}
    
    \renewcommand{\arraystretch}{1.5}
    \setlength{\tabcolsep}{8pt}
    
    \begin{tabularx}{\textwidth}{
        >{\centering\arraybackslash}p{0.12\textwidth}
        >{\centering\arraybackslash}p{0.16\textwidth}
        >{\raggedright\arraybackslash}X
        >{\centering\arraybackslash}p{0.16\textwidth}
    }
        \toprule
        \textbf{参数}
        & \textbf{Spearman $\rho$}
        & \textbf{观察到的边际关系}
        & \textbf{结论强度} \\
        \midrule

        $C_1$
        & 0.058
        & 几乎无稳定单调关系
        & 弱 \\

        $Q_1$
        & -0.886
        & 较强负向关系：当前设计点中 $Q_1$ 越高，早期退化率总体越低
        & 中等方向性 \\

        $C_2$
        & 0.406
        & 弱至中等正向关系：较高 $C_2$ 倾向对应较高退化
        & 弱 \\

        三参数综合
        & ---
        & 参数高度耦合，无法把上述边际关系直接解释为独立因果作用
        & --- \\

        \bottomrule
    \end{tabularx}
\end{table}

$C_1,Q_1,C_2$分别得到$\rho=0.058$（$p=0.9333$）、$\rho=-0.886$（$p=0.0333$）和$\rho=0.406$（$p=0.4333$）。由数据可知，在当前六个独立策略设计点中，$Q_1$与早期退化率表现出最强的负向边际单调关系，即阶段切换SOC较高的策略总体倾向于具有较低的早期退化率；$C_2$仅表现出较弱的正向趋势，而$C_1$基本未观察到稳定单调关系。从边际关联强度看，$Q_1$明显强于另外两个参数。但由于三参数Holm校正后$Q_1$的$p=0.100$，且各参数并非正交设计，因此上述关系只能解释为\textbf{当前策略设计下的边际关联方向}，不能解释为独立因果效应。

\begin{figure}[htbp]
\centering
\includegraphics[width=0.62\textwidth]{texfile/figures/fig_p2_q1_exact.pdf}
\caption{$Q_1$与策略中位主退化率的精确相关}
\label{fig:p2q1}
\end{figure}

\subsubsection{参数联合可识别性与影响程度}
 现有数据不足以可靠量化 $C_1,Q_1,C_2$ 各自的独立影响程度。原因是六个独立策略点中参数设计具有明显共线性，尤其 $C_1$ 与 $C_2$ 的相关系数约为 $-0.823$；进一步建立的策略级Ridge模型在整策略留出验证中的 $R^2$ 仅为0.046，说明联合模型对未见策略缺乏可靠泛化能力。因此，\textbf{不能根据Ridge系数给出可信的独立影响大小排序}。若仅比较充电参数的边际关联强度，则有
\[
|\rho_{Q_1}|>|\rho_{C_2}|>|\rho_{C_1}|,
\]
即 $Q_1$ 是当前数据中最值得优先关注的参数，$C_2$ 次之，$C_1$ 最弱；但这一排序只是\textbf{边际关联强度排序}，\textbf{不是独立影响程度或因果重要性排序}。正则化参数、系数和删变量结果见支撑材料。

\subsubsection{不同SOC区间的倍率暴露效应}
为区分倍率大小与施加位置，将策略$p$表示为SOC上的实际倍率函数，并在四个预设区间计算平均暴露：
\begin{equation}
C_p(s)=\begin{cases}C_{1,p},&0\le s<Q_{1,p},\\C_{2,p},&Q_{1,p}\le s\le80,
\end{cases}\qquad
E_{p,k}=\frac{1}{|I_k|}\int_{I_k}C_p(s)\,\mathrm ds,
\end{equation}
其中$I_k\in\{[0,20),[20,40),[40,60),[60,80]\}$。在短区间容量和内阻近似稳定时，$I^2R\,\mathrm dt\propto C\,\mathrm ds$，故$E_{p,k}$可作为单位SOC区间的近似倍率—欧姆热暴露。S8在两个中高SOC区间均具有最高暴露，见图\ref{fig:p2exposure}。

\begin{figure}[htbp]
\centering
\includegraphics[width=0.64\textwidth]{texfile/figures/fig_p2_soc_exposure.pdf}
\caption{S4--S9在四个SOC区间的平均倍率暴露}
\label{fig:p2exposure}
\end{figure}

进一步定义
\begin{equation}
E_{L,p}=\frac{E_{p,1}+E_{p,2}}{2},\qquad E_{H,p}=\frac{E_{p,3}+E_{p,4}}{2},\qquad
E_{\mathrm{mean},p}=\frac{E_{L,p}+E_{H,p}}{2},\qquad E_{\mathrm{contrast},p}=E_{H,p}-E_{L,p}.
\label{eq:p2-low-high-exposure}
\end{equation}
$E_{\mathrm{mean},p}$与退化率呈正向但不显著（$\rho=0.600,p=0.2417$），$E_{\mathrm{contrast},p}$无稳定单调证据（$\rho=0.143,p=0.8028$）。替代SOC切分下中高SOC方向仍为正但均未显著，因此只保留方向性弱证据。局部准匹配结果见表\ref{tab:p2localpairs}；其中S8$\to$S9伴随中高SOC倍率降低时两项退化指标同步下降，但三个参数同时变化，不能解释为单参数因果效应。综上所述，本文将“中高SOC高倍率可能增加退化风险”保留为方向性弱证据，并将其用于后续局部策略建模，而不视为已确认的独立SOC区间效应。

\begin{table}[htbp]
\centering
\caption{局部准匹配策略对比}
\label{tab:p2localpairs}
\footnotesize
\begin{tabularx}{\textwidth}{p{1.5cm}p{2.5cm}cccXX}
\toprule
对比 & 共享条件 & $\Delta C_1$ & $\Delta Q_1$ & $\Delta C_2$ & $\Delta r$（百分点/100次） & $\Delta D$/\%\\
\midrule
S4$\to$S5 & $C_2=4.0$C & $+0.3$C & $-13$ & 0 & $-0.0146$ & $-0.0066$\\
S6$\to$S7 & $C_1=5.6$C & 0 & $+17$ & $-0.3$C & $-0.0091$ & $-0.0539$\\
S8$\to$S9 & 同组、同结构标签 & $+1.1$C & $+49$ & $-1.1$C & $-1.7735$ & $-2.1470$\\
\bottomrule
\end{tabularx}
\end{table}

\subsubsection{退化机理一致性与证据边界}
在基础SOC暴露模型上分别加入Cycle 80--150平均温度和内阻相对变化。加入温度后整策略留出RMSE由0.636数值下降至0.347，仅加入内阻未产生类似改善；温度与策略中位退化率的精确Spearman为$\rho=0.657,p=0.1750$。温度是在策略实施后观测的伴随变量，其预测增益与高倍率热效应、极化和界面副反应的机理方向一致\cite{koehler2019,operando2022}，但不能解释为已验证的因果中介；内阻在当前早期窗口也未表现出稳定增益。

综上，问题二形成分层证据：总体策略差异及S8高退化属于正式强证据；$Q_1$负向关系是方向稳定的中等证据；Ridge与SOC区间模型仅提供可识别性和方向性探索；温度、内阻承担机理辅助。现有数据不能对三个快充参数建立可靠的独立因果重要性排序。问题四继续沿用$r_i$并聚合为六个策略响应点，SOC暴露只作为连续代理的预测坐标。
\FloatBarrier

\subsection{问题三：嵌套验证的短期SOH预测与多模型EOL外推}
\subsubsection{预测边界与特征体系}
9块测试电池仅观测至第150次循环，40块非测试电池的第151--200次记录可用于回测；全部电池均未达到80\% SOH。因此，将任务分为可验证的第151--200次短期预测和无真实标签的80\% EOL条件外推，二者采用独立的数据流与评价口径。

为保持与问题一、二一致，稳定基准和正式相对SOH始终沿用式\eqref{eq:relative-soh}：
\[
B_i=\operatorname{median}_{20\le n\le30}S^c_{i,n},\qquad
R_{i,n}=\frac{S^c_{i,n}}{B_i}.
\]
趋势特征和经验曲线均基于清洗序列。令$\ell\in\{50,100,150\}$为可见早期长度，$\boldsymbol x_i^{(\ell)}$包括健康水平、退化动态、物理策略、实验元数据和运行信号。物理变量与\texttt{dataset\_id}等元数据分组处理，全部插补、标准化、特征筛选和降维仅在训练折拟合。

对单段3.6C策略S2，原始$C_1$缺失值仍保持\texttt{NaN}；模型特征层另设单段协议标识，并按0--80\% SOC全程3.6C构造有效倍率与SOC暴露特征，从而避免将统计插补值误解释为真实物理倍率。

\begin{figure}[htbp]
\centering
\includegraphics[width=0.67\textwidth]{texfile/figures/fig01_revised_flow.png}
\caption{问题三的信息边界与平衡嵌套验证流程}
\label{fig:p3flowv2}
\end{figure}

\subsubsection{策略约束的嵌套验证与短期Hybrid模型}
以电池为独立样本，外层采用2轮$\times$7折的策略约束Round-robin验证，共形成14个Outer折。每轮将各策略内电池轮转分配至不同验证折，使40块非测试电池在每轮均恰好进入验证1次，两轮后每块电池总计进入Outer验证2次；训练端内部仍采用4折策略平衡Inner CV。所有窗口搜索、插补、标准化、特征筛选、PCA、候选模型及混合权重均仅在Outer训练端内部确定。窗口搜索仅用于预测特征工程：
\begin{equation}
\begin{aligned}
r_{i,\mathrm{pred}}^{(a)}&=-\operatorname{median}_{a\le u<v\le150}
\frac{R_{i,v}-R_{i,u}}{v-u},\\
r_{i,\mathrm{fut}}&=-\operatorname{median}_{150\le u<v\le200}
\frac{R_{i,v}-R_{i,u}}{v-u},\\
a^*&=\arg\max_{a\in\{60,70,80,90,100\}}
\rho_S\!\left(r_{i,\mathrm{pred}}^{(a)},r_{i,\mathrm{fut}}\right).
\end{aligned}
\label{eq:p3-window}
\end{equation}
每个外层折只用训练电池的后续段计算$r_{i,\mathrm{fut}}$，验证电池的第151--200次不参与选择。全部40块训练电池定型时仍得到$a^*=100$；14个Outer折中，起点90和100次分别出现2次和12次，未再选择60、70或80作为最优起点，说明后段趋势特征在平衡验证下更加稳定。

对第$i$块电池，令未来50维轨迹相对第150次的变化为
\begin{equation}
\Delta\mathbf y_i=
\left(S^c_{i,151}-S^c_{i,150},\ldots,S^c_{i,200}-S^c_{i,150}\right)^\top.
\label{eq:p3-delta-y}
\end{equation}
在Inner训练折中对$\Delta\mathbf y_i$做PCA并保留累计解释方差不低于99\%的主成分，再分别拟合PCA--Ridge、PCA--GPR和PCA--SVR；经验候选为拉伸指数、幂律损失和非负二次损失。Inner CV分别选择最优数据模型和经验模型，并仅在Inner预测上选择混合权重：
\begin{equation}
\widehat{\mathbf S}^{\rm H}_{i,151:200}=w\widehat{\mathbf S}^{\rm D}_{i,151:200}+
(1-w)\widehat{\mathbf S}^{\rm E}_{i,151:200},\qquad 0\le w\le1.
\label{eq:p3-hybrid}
\end{equation}
全训练集定型后的数据模型为PCA--Ridge、经验模型为拉伸指数，$w=0.35$。这表示数据驱动分量承担35\%的轨迹校正，而不是由黑箱模型取代受约束的退化趋势。

需要强调的是，上述固定组合仅用于全部40块训练电池完成方法选择后的最终测试集预测；在14个Outer折中，数据模型、经验模型和$w$均由该折Inner CV独立选择，因此下述Outer误差评价的是完整嵌套选择流程，而非固定PCA--Ridge与拉伸指数组合的性能。在14个平衡外层折上计算各模型的 RMSE，结果如图\ref{fig:p3nested}所示：

\begin{figure}[htbp]
\centering
\includegraphics[width=0.68\textwidth]{texfile/figures/fig02_nested_model_comparison.png}
\caption{14个平衡Outer折上的短期模型比较}
\label{fig:p3nested}
\end{figure}

在14个平衡Outer折上，嵌套选择式Hybrid的折级平均RMSE为0.000480 SOH（0.0480个百分点，标准差0.000312），低于Inner选择的最佳单模型0.000509；平均MAE为0.000278，Cycle 200绝对误差为0.000473，相对最佳单模型的平均RMSE改善5.7\%。40块训练电池均恰好进入Outer验证2次；按电池汇总两轮验证误差后，30/40块（75.0\%）的Hybrid轨迹RMSE低于对应最佳单模型。按策略等权的macro-policy RMSE为0.000506。该结果反映完整嵌套选择流程的泛化误差，而非固定最终组合的外层性能。

\subsubsection{预测区间、早期数据长度与测试电池预测}
对第$o$个外层训练集，仅使用其Inner OOF电池的50步轨迹最大绝对残差，取95\%分位数$q_{0.95}^{(o)}$构造该折外层验证区间：
\begin{equation}
\mathcal I_{i,n}^{(o)}=\left[\hat S_{i,n}-q_{0.95}^{(o)},\ \hat S_{i,n}+q_{0.95}^{(o)}\right],
\qquad n=151,\ldots,200.
\label{eq:p3-pi}
\end{equation}
除点级覆盖率PICP外，再定义整轨迹覆盖率
\begin{equation}
\mathrm{TICP}=\frac{1}{\sum_o|\mathcal V_o|}\sum_o\sum_{i\in\mathcal V_o}
\mathbf 1\!\left[\max_{151\le n\le200}|S^c_{i,n}-\hat S_{i,n}|\le q_{0.95}^{(o)}\right],
\label{eq:p3-ticp}
\end{equation}
其中$\mathcal V_o$为第$o$次Outer-validation电池集。按全部Outer验证轨迹加权统计，PICP为98.8\%、TICP为92.5\%，平均区间宽度MPIW为0.005311 SOH（0.5311个百分点）。点级覆盖仍略偏保守，且整条50点轨迹同时覆盖更严格，不能由PICP替代。

为公平评价早期观测长度，固定PCA--Ridge架构并使$\ell=50,100,150$均预测同一第151--200次目标；每种长度只使用自身可见区间内的特征。结果见表\ref{tab:p3length}：从50增加到100次、从100增加到150次，RMSE分别下降54.6\%和56.0\%，因此，早期循环信息并未在前50或100次后饱和，增加可观测循环长度能够显著提高后续SOH预测精度；在本数据范围内，使用完整前150次数据具有明显优势。

\begin{table}[htbp]
\centering
\caption{不同早期观测长度的嵌套验证结果（统一预测Cycle 151--200）}
\label{tab:p3length}
\begin{tabular}{ccccc}
\toprule
$\ell$ & RMSE & RMSE标准差 & MAE & Cycle 200绝对误差\\
\midrule
50  & 0.002842 & 0.002076 & 0.001795 & 0.002170\\
100 & 0.001291 & 0.000457 & 0.000819 & 0.001093\\
150 & 0.000568 & 0.000379 & 0.000283 & 0.000435\\
\bottomrule
\end{tabular}
\end{table}

用全部40块训练电池重新完成训练端选择后，对9块测试电池一次性预测第151--200次SOH，见图\ref{fig:p3testv2}。对9块测试电池预测发现，各电池在第151–200次总体仍保持缓慢衰减，其中电池16（S8）第200次预测SOH最低，为96.259\%；电池9（S3）为98.149\%，其余多数电池仍在99.2\%左右及以上。因此，S8测试电池表现出最明显的短期容量衰减，其余测试电池在第200次循环前总体仍保持较高SOH。
该排序与问题一的早期衰减表征方向一致，但它是预测输出，不能反向充当问题二显著性检验的证据。
\begin{figure}[htbp]
\centering
\includegraphics[width=0.82\textwidth]{texfile/figures/fig07_test_predictions.png}
\caption{九块测试电池第151--200次SOH预测}
\label{fig:p3testv2}
\end{figure}

\subsubsection{80\% EOL的多模型约束外推}
由于没有电池到达80\%阈值，EOL不使用短期Hybrid预测再拟合。测试电池仅用真实Cycle 80--150观测，训练电池则比较80--150与80--200拟合参数，以校准早期拟合偏差。

设$t=n-30$，在相对基准SOH尺度上保留三种单调退化模型族：
\begin{equation}
\begin{aligned}
\mathcal M_{\rm SE}:&\quad \hat R(n)=\exp(-at^b), &&a\ge\varepsilon,\ 0.2\le b\le4,\\
\mathcal M_{\rm Pow}:&\quad \hat R(n)=1-at^b, &&a\ge\varepsilon,\ 0.2\le b\le4,\\
\mathcal M_{\rm Quad}:&\quad \hat R(n)=1-at-ct^2, &&a,c\ge\varepsilon,
\end{aligned}
\qquad \varepsilon=10^{-12}.
\label{eq:p3-eol-families}
\end{equation}
严格正下界保证所有待校准参数的对数均有定义。令$\boldsymbol\theta_{i,u:v}^{(m)}$为电池$i$在窗口$u{:}v$上拟合的模型$m$参数向量，并定义同一训练电池的联合校准向量
\begin{equation}
\boldsymbol\eta_i^{(m)}=\log\boldsymbol\theta_{i,80:200}^{(m)}-
\log\boldsymbol\theta_{i,80:150}^{(m)}.
\label{eq:p3-eta}
\end{equation}
对策略$p$的校准均值采用向全局均值收缩的经验贝叶斯形式
\begin{equation}
\boldsymbol\mu_p^{(m)}(\kappa)=\frac{n_p}{n_p+\kappa}\overline{\boldsymbol\eta}_p^{(m)}+
\frac{\kappa}{n_p+\kappa}\overline{\boldsymbol\eta}_{\rm global}^{(m)},\qquad \kappa=3.
\label{eq:p3-eb}
\end{equation}
Monte Carlo每次从同一训练电池$j$联合抽取完整向量，而不独立拼接参数分量：
\begin{equation}
\boldsymbol\eta^{*(m)}=\boldsymbol\mu_p^{(m)}+
\left(\boldsymbol\eta_j^{(m)}-\overline{\boldsymbol\eta}_{\rm global}^{(m)}\right),\qquad
\boldsymbol\theta_{i,\rm corr}^{(m)}=\boldsymbol\theta_{i,80:150}^{(m)}\odot
\exp\!\left(\boldsymbol\eta^{*(m)}\right).
\label{eq:p3-joint-calibration}
\end{equation}
联合抽样保留$(a,b)$或$(a,c)$的经验相关结构，并使校准中心服从策略层收缩。每个模型族独立求解$\hat R(n)\le0.8/B_i$的首次$n$，得到EOL样本$N_{i,\rm EOL}^{(m)}$。跨模型外推包络定义为
\begin{equation}
L_{i,\rm EOL}=\min_m Q_{0.025}\!\left(N_{i,\rm EOL}^{(m)}\right),\qquad
U_{i,\rm EOL}=\max_m Q_{0.975}\!\left(N_{i,\rm EOL}^{(m)}\right).
\label{eq:p3-envelope}
\end{equation}
以 SOH 降至初始值的 80\% 作为寿命终点，本文基于预测模型估计各电池达到该阈值时对应的循环次数，并采用 Bootstrap 重采样构建 95\% 预测区间。如图\ref{fig:p3eolenv}所示；不同电池的 80\% SOH 寿命预测存在明显差异。圆点表示寿命点预测值，水平线表示 2.5\%–97.5\% Bootstrap 区间。S9 的预测寿命最高，而 S3 与 S8 的预测寿命相对较短；部分电池的区间较宽，表明其寿命外推的不确定性较高。因此，寿命预测结果应结合区间范围进行解释，而不能仅依据点预测值进行排序。

\begin{figure}[htbp]
\centering
% 左侧：图片
\begin{minipage}[c]{0.48\textwidth}
\centering
\includegraphics[width=\linewidth]{texfile/figures/fig09_eol_envelope_log.png}
\caption{测试电池80\% EOL的跨模型外推包络}
\label{fig:p3eolenv}
\end{minipage}
\hfill
% 右侧：表格
\begin{minipage}[c]{0.48\textwidth}
\centering
\small
\begin{tabular}{ccrrrr}
\toprule
电池$i$ & 策略$p(i)$ & $\widetilde N_{i,\rm EOL}$ & 包络下界 & 包络上界 & $R_{\rm ext,i}$\\
\midrule
9  & S3 & 798 & 202 & 1,279 & 12.0\\
16 & S8 & 1,167 & 206 & 3,537 & 19.3\\
24 & S7 & 1,376 & 495 & 2,530 & 23.5\\
5  & S2 & 1,453 & 510 & 2,338 & 25.1\\
14 & S6 & 1,512 & 355 & 2,249 & 26.2\\
10 & S4 & 3,367 & 647 & 4,706 & 63.3\\
2  & S1 & 4,606 & 411 & 6,752 & 88.1\\
11 & S5 & 6,870 & 957 & 9,300 & 133.4\\
25 & S9 & 14,895 & 233 & 27,552 & 293.9\\
\bottomrule
\end{tabular}
\caption{测试电池80\% EOL的多模型条件外推结果}
\label{tab:p3eol}
\end{minipage}
\end{figure}

令$\mathcal M=\{\mathcal M_{\rm SE},\mathcal M_{\rm Pow},\mathcal M_{\rm Quad}\}$。多模型中位中心及电池级外推指数统一定义为
\begin{equation}
\widetilde N_{i,\rm EOL}=\operatorname{median}_{m\in\mathcal M}
\left\{Q_{0.5}\!\left(N_{i,\rm EOL}^{(m)}\right)\right\},\qquad
R_{\rm ext,i}=\frac{\widetilde N_{i,\rm EOL}-200}{50}.
\label{eq:p3-eol-center}
\end{equation}
由于所有电池均未实际达到80\% SOH，本文得到的EOL只能解释为模型条件下的寿命外推，而不能视为已观测真实寿命。从多模型中位中心看，9块测试电池由短到长大致排序为
\[
S3 < S8 < S7 < S2 < S6 < S4 < S1 < S5 < S9.
\]
其中电池9（S3）的EOL中位中心约为798次，电池16（S8）约为1167次；电池11（S5）约为6870次，电池25（S9）约为14895次，表现出较长寿命倾向。但远期外推区间较宽，例如S9的跨模型包络为233--27552次、外推指数达到293.9，因此\textbf{EOL中心主要用于比较不同电池的相对寿命倾向，不能把14895次等数值解释为已经验证的实际循环寿命}。


\subsubsection{稳健性与特征解释}
对收缩强度、拟合起点和指数上界进行敏感性分析。收缩强度变化时EOL排序相关不低于0.983，中心值中位相对变化不超过2.94\%；起点改为70和90时，排序相关分别为0.983和0.933，中心值变化6.68\%和15.36\%。相对排序较稳定，但绝对EOL仍对窗口敏感。

在同一平衡Outer划分下，逐组消融显示，依次加入退化动态、物理策略、元数据和运行信号后，PCA--Ridge数据驱动分支的平均RMSE由0.001413降至0.000946；其中物理策略带来的直接增益较小，而元数据与运行信号仍提供一定预测增量。组置换进一步表明，在PCA--Ridge数据驱动分支中，健康水平是最主要的信息源，退化动态次之，元数据与运行信号提供较小增量，物理策略的直接置换增量最小，如图\ref{fig:p3permutation}所示。该结果仅解释数据驱动分支的预测依赖，不代表整个Hybrid模型的因果效应。
\begin{figure}[htbp]
\centering
\includegraphics[width=0.66\textwidth]{texfile/figures/fig12_group_permutation.png}
\caption{PCA--Ridge数据驱动分支的特征组置换重要性}
\label{fig:p3permutation}
\end{figure}
\FloatBarrier


\subsection{问题四：基于前三问结论的风险感知双目标优化}
\subsubsection{前述证据衔接与策略级响应构建}
问题一确定Cycle 80--150退化率$r_i$及典型策略，问题二确认参数推断的独立单位为S4--S9六个策略点，问题三表明80\% EOL只能辅助判断寿命方向。基于这些结论，将40块重复电池聚合为策略级代表退化率
\begin{equation}
r_{p,\mathrm{rep}}=\mathop{\mathrm{median}}_{i\in p}r_i,\qquad
Y_p=\ln\!\left(\frac{r_{p,\mathrm{rep}}}{r_0}\right),\quad r_0=10^{-5},
\label{eq:p4-health-response}
\end{equation}
并以策略级实测充电时间中位数$T_p$（取问题一$\tau_p^{\rm cycle}$的实测时间口径）与$Y_p$为两个越小越好的响应。

S4--S9的$r_{p,\mathrm{rep}}$与问题三多模型EOL中位中心在模型内部呈较强反向排序一致性（$\rho=-0.943$，720种排列双侧$p_{\rm exact}=0.0167$）。由于EOL中心本身由早期SOH信息及训练电池校准构造，该相关性仅作为模型内部一致性检查，不视为独立远期寿命验证；因此问题四仍以可直接观测的早期退化率作为健康目标。因此，问题四不重复建立已有策略的退化指标，而直接继承问题一经后续循环验证的Cycle 80--150稳健退化率，并在策略层取重复电池中位数作为各快充策略的SOH衰减代表值。在S4--S9六种策略中，早期SOH衰减率由低到高依次为
\[
S9 < S5 < S4 < S7 < S6 \ll S8,
\]
对应策略中位退化率分别为
\[
1.997,\; 2.676,\; 2.822,\; 2.990,\; 3.081,\; 19.732 \times 10^{-5}/\text{cycle}.
\]
因此，S9在现有实验策略中SOH衰减最缓，S8衰减最严重，S4--S7中的S4、S5、S6、S7总体处于相近的中等退化水平。上述实测策略级退化结果作为问题四健康目标的基础数据；后续健康Guard主要用于估计尚未实际实验的连续充电策略，而不是重新定义已有策略的SOH衰减。

\subsubsection{充电时间代理与 Guard}
令连续策略为$\boldsymbol x=(C_1,Q_1,C_2)$。两阶段恒流的电量守恒基线为
\begin{equation}
T_{\rm CC}(\boldsymbol x)=0.6\left(\frac{Q_1}{C_1}+\frac{80-Q_1}{C_2}\right)\quad(\mathrm{min}).
\label{eq:p4-tcc}
\end{equation}

由两阶段恒流充电时间模型可知，在其他参数不变时，第一阶段和第二阶段充电倍率 $C_1,C_2$ 增大均会单调缩短充电时间。阶段切换SOC $Q_1$ 的影响则取决于两个阶段倍率的相对大小：当 $C_1 > C_2$ 时，提高 $Q_1$ 相当于扩大较高倍率第一阶段的SOC覆盖范围，因此充电时间缩短；当 $C_1 < C_2$ 时结论相反；当 $C_1 = C_2$ 时，理论充电时间与 $Q_1$ 无关。因此，\textbf{充电时间不仅取决于倍率大小，还取决于高倍率被分配在哪一段SOC区间以及该区间的宽度。}

在每个训练折内，一条分支以训练策略时间残差$e_T=T_p-T_{\rm CC}$的中位数校正物理基线，另一条用Matern核高斯过程拟合残差：
\begin{equation}
T_{\rm base}=T_{\rm CC}+\mathop{\mathrm{median}}(e_T),\qquad
T_{\rm GP}=T_{\rm CC}+\widehat e_{T,\rm GP},\qquad
\mu_T=\max\{T_{\rm base},T_{\rm GP}\}.
\label{eq:p4-time-guard}
\end{equation}
取较慢分支以抑制系统性低估。时间Guard采用策略级Outer LOPO，GP核参数只在训练端选择；其RMSE为0.4379 min、MAE为0.2975 min、$R^2=-0.419$，过度乐观率为16.7\%。据此定义
\begin{equation}
\sigma_T(\boldsymbol x)=\max\{s_{T,\rm guard},\sigma_{T,\rm GP}(\boldsymbol x)\},
\qquad s_{T,\rm guard}=0.4379\ \mathrm{min}.
\label{eq:p4-time-sigma}
\end{equation}

\subsubsection{健康退化代理与 Guard}
沿用问题二预先固定的四个SOC区间倍率暴露$E_{p,k}$，并记低、高SOC两区间平均暴露为$E_{L,p},E_{H,p}$。低复杂度分支采用
\begin{equation}
\mu_{\rm Ridge}(\boldsymbol x)=\beta_0+\beta_L Z(E_L(\boldsymbol x))+\beta_H Z(E_H(\boldsymbol x)),
\end{equation}
局部保护分支以四区间暴露为输入建立Matern-GP。为防止低估退化，最终健康均值与不确定性取
\begin{equation}
\mu_Y=\max\{\mu_{\rm Ridge},\mu_{\rm GP}\},\qquad
\sigma_Y=\max\{\sigma_{\rm Ridge},\sigma_{\rm GP},s_{Y,\rm guard}\}.
\label{eq:p4-health-guard}
\end{equation}
Ridge正则系数和GP超参数均在Outer训练端选择，$s_{Y,\rm guard}$由max-Guard外层残差标定。健康Guard的RMSE为0.9060 log-rate、MAE为0.6090 log-rate、$R^2=-0.432$，过度乐观率为16.7\%。两个负$R^2$说明六个策略点不足以支持高精度整策略外推，连续模型只能生成待验证候选。

\begin{figure}[!htbp]
\centering
\includegraphics[width=0.94\textwidth]{texfile/figures/fig14_guard_outer_lopo_v2.png}
\caption{时间与健康Guard的Outer LOPO验证}
\label{fig:p4guardlopo}
\end{figure}

\subsubsection{可信域约束与风险 Pareto 优化}
为避免把三个参数各自的边界自由拼成无实验支撑的角点，连续搜索域同时施加凸包和最近邻约束。令$Z(\cdot)$为按六个策略点拟合的参数标准化，定义
\begin{equation}
\mathcal X_{\rm hull}=\operatorname{Conv}\{\boldsymbol x_{\rm S4},\ldots,\boldsymbol x_{\rm S9}\},\qquad
d_{\min}(\boldsymbol x)=\min_{p\in\{\rm S4,\ldots,S9\}}\|Z(\boldsymbol x)-Z(\boldsymbol x_p)\|_2,
\end{equation}
\begin{equation}
\mathcal X_{\rm safe}=\{\boldsymbol x\in\mathcal X_{\rm hull}:d_{\min}(\boldsymbol x)\le\delta_0\}.
\label{eq:p4-safe-domain}
\end{equation}
阈值$\delta_0$取实验策略最近邻距离的90\%分位，本数据中为2.367；Sobol低差异凸组合和边界插值共保留71542个可信域候选点（受限搜索点）。

在$\mathcal X_{\rm safe}$内建立风险修正目标
\begin{equation}
F_T(\boldsymbol x)=\mu_T(\boldsymbol x)+\kappa_T\sigma_T(\boldsymbol x),\qquad
F_Y(\boldsymbol x)=\mu_Y(\boldsymbol x)+\kappa_Y\sigma_Y(\boldsymbol x),
\label{eq:p4-risk-objectives}
\end{equation}
基准情景取$\kappa_T=\kappa_Y=1$；$\kappa$表示风险厌恶尺度，不解释为严格概率分位。对Pareto前沿内两个目标作极差归一化，建立风险修正双目标优化模型
\begin{equation}
\boldsymbol x_{\rm Bal}=\arg\min_{\boldsymbol x\in\mathcal P}
\sqrt{\widetilde F_T(\boldsymbol x)^2+\widetilde F_Y(\boldsymbol x)^2}.
\label{eq:p4-balanced}
\end{equation}
基于上述可信域约束与风险修正双目标，得到连续候选策略的可行分布及其 Pareto 前沿，如图\ref{fig:p4trustpareto}所示。
\begin{figure}[!htbp]
\centering
\includegraphics[width=0.94\textwidth]{texfile/figures/fig15_trust_pareto_v2.png}
\caption{可信域及连续风险Pareto前沿}
\label{fig:p4trustpareto}
\end{figure}

\subsubsection{Pareto 求解与分级推荐}
由式（30），在可信域内首先依据风险修正后的充电时间目标 $F_T$ 和SOH衰减目标 $F_Y$ 筛选连续Pareto前沿。由于Pareto前沿包含多个互不支配方案，为获得兼顾充电速度与健康损失的代表性折中策略，本文分别对两项目标进行极差归一化，并选取距离理想点 $(0,0)$ 最近的Pareto方案作为Balanced候选。求解得到Balanced为
\[
(C_1, Q_1, C_2) = (5.480C,\, 22.39\%,\, 4.656C),
\]
预测充电时间约为9.935 min，预测早期退化率约为 $3.617 \times 10^{-5}/\text{cycle}$。

为判断连续Balanced候选是否相较已有实验策略具有实际优势，首先不依赖连续代理模型，而直接以S4--S9的实测充电时间 $T_p$ 和策略中位早期退化率 $r_{p,\text{rep}}$ 进行非支配筛选。按未取整数据计算时，S7、S5和S9均位于数学Pareto前沿；但S7相较S5的充电时间仅缩短0.000147 min（约0.0089 s），低于本文0.001 min的时间报告精度，因此这一差异不足以构成具有实际意义的速度优势。采用 $\varepsilon_T=0.001$ min 的实践Pareto准则后，已有实测策略的有效前沿归并为S5和S9。进一步通过20000次策略内电池级Bootstrap检验前沿稳定性，S5和S9进入经验Pareto前沿的频率分别为63.0\%和83.7\%；同时S4和S7的前沿频率分别为51.4\%和59.4\%。因此，S5更适合作为已有策略中速度/均衡区域的代表，而不是唯一统计最优策略；S9则具有当前最稳定的健康优先端证据。表\ref{tab:p4-recommendations}列出了已有实测策略 S5、S9 与连续候选策略 Fast、Balanced、Life 的参数配置、充电时间及预测主退化率，为后续 Pareto 优势检验提供对照依据。
\begin{table}[H]
\centering
\caption{已有策略Pareto前沿与连续代表候选}
\label{tab:p4-recommendations}
\small
\begin{tabular}{lcccccc}
\toprule
类型 & 方案 & $C_1$/C & $Q_1$/\% & $C_2$/C & 时间/min & $r\,(10^{-5}/\mathrm{cycle})$\\
\midrule
已有实测 & S5 & 5.300 & 54.00 & 4.000 & 10.043 & 2.676\\
已有实测 & S9 & 4.800 & 80.00 & 4.800 & 11.038 & 1.997\\
连续候选 & Fast & 4.638 & 25.08 & 5.258 & 9.701 & 15.483\\
连续候选 & Balanced & 5.480 & 22.39 & 4.656 & 9.935 & 3.617\\
连续候选 & Life & 4.799 & 50.29 & 4.759 & 10.365 & 2.001\\
\bottomrule
\end{tabular}
\end{table}

连续前沿含332个非支配点。Balanced相对S5的预测时间改善约0.108 min，且预测退化率增幅约35.1\%，但是明显小于时间Guard的Outer-LOPO RMSE 0.438 min，因此现模型不足以证明Balanced真实比S5更快，更不能认为其综合性能优于S5。Balanced 相对 S9的预测时间改善约 1.103min，退化率约提高81.1\%，因此，Balanced虽然在模型意义下获得更快的充电速度，但明显牺牲了健康目标，不能替代S9作为健康优先策略。此外，健康Guard的整策略Outer‑LOPO $R^2=-0.432$，进一步说明连续候选的退化预测仍存在较强外推不确定性。
结合已有策略的实测Pareto结果、Bootstrap稳定性以及连续模型的外推误差，本文不将Balanced直接视为最终最优策略，而形成如下分级推荐：若优先考虑充电速度并兼顾较低SOH衰减，推荐S5（5.3C–54\%–4.0C）作为速度/均衡场景代表；若优先考虑降低SOH衰减和保持电池健康，则推荐S9（4.8C–80\%–4.8C）作为健康优先策略。 Balanced（5.480C–22.39\%–4.656C）虽然是风险Pareto前沿上两项目标归一化后距离理想点最近的连续折中解，但其相较S5的预测时间改善小于时间模型的外层预测误差，且预测退化率更高，因此仅作为下一轮独立老化实验的优先验证候选，在获得新的实测证据前不建议替代S5或S9。
\subsubsection{稳健性与适用范围}
将可信距离阈值取$0.75\delta_0,\delta_0,1.25\delta_0$时，可信域候选点数分别为69861、71542、71542，Balanced均保持在同一候选点。风险情景取$(\kappa_T,\kappa_Y)=(0,0),(1,1),(0.5,1.96),(1.96,1.96)$时，Balanced也不漂移；原因是当前可信域内Outer-LOPO误差下限高于局部分支标准差，风险项主要增加共同误差基线。这一现象不能解释为模型已消除风险，反而说明主要不确定性来自整策略泛化。

对连续前沿作1200次不确定性抽样，Fast、Balanced和Life区域出现非支配点的概率分别为58.4\%、95.7\%和90.2\%。此外，20,000次策略内电池级Bootstrap中，S4、S5、S6、S7、S8、S9进入经验Pareto的频率分别为51.4\%、63.0\%、15.2\%、59.4\%、0\%和83.7\%，说明速度端存在明显重采样不确定性，而S9健康端相对更稳定。因此连续解只表述为候选区域；S5作为速度/均衡簇代表而非唯一最优，S9的健康优先推荐不依赖连续代理。

推荐范围限于原始分组3的新结构实验层或高度相近的A123 18650磷酸铁锂电芯、0--80\% SOC两阶段恒流协议及相同温控和放电条件。超出$\mathcal X_{\rm safe}$、改变材料体系、容量、温度边界或加入CV阶段时，均须重新标定代理、可信域和不确定性。
\FloatBarrier


\section{模型检验、误差与敏感性分析}
\subsection{独立性与信息泄漏检查}
策略间分布比较以独立电池为样本；涉及$C_1,Q_1,C_2$的参数作用推断及问题四连续代理均以策略为独立设计单位；问题三按\texttt{battery\_id}整体分组，始终不把逐循环记录视为独立样本。问题三的窗口、插补、特征筛选、PCA、候选模型与Hybrid权重均限于相应Outer训练端，且特征只使用截止时刻以前的数据，从而避免循环级随机划分、未来信息和模型选择泄漏。该边界与按电芯批次报告预测误差的做法一致\cite{wang2021}。

\subsection{问题一、二的稳健性}
问题一、二的主要结论对指标和窗口选择具有方向稳定性，汇总见表\ref{tab:q12-robustness}。其中替代响应和多种扰动用于检验结论敏感性，并不构成彼此独立的重复证据。

\begin{table}[htbp]
    \centering
    \caption{问题一、二稳健性总览}
    \label{tab:robustness_summary}
    \small
    \renewcommand{\arraystretch}{1.25}
    \setlength{\tabcolsep}{6pt}

    \begin{tabularx}{\textwidth}{
        >{\raggedright\arraybackslash}p{0.18\textwidth}
        >{\raggedright\arraybackslash}l
        >{\raggedright\arraybackslash}X
    }
        \toprule
        检验对象 & 扰动方式 & 稳健性结论 \\
        \midrule

        基准与退化窗口
        & 更换稳定基准及 Cycle 60--150 至 100--150
        & S9/S1 持续位于低退化端，S8 持续最高 \\

        替代指标
        & $r_i$ 与 $D_i$ 互换
        & 典型 S9/S8 判断一致，但不视为独立验证 \\

        总体策略差异
        & $D_i$ 复核及逐策略删除
        & $D_i$ 全局 $p=0.0111$；仅删 S8 后 $p=0.1173$ \\

        $Q_1$ 方向
        & 不同窗口、$D_i$ 及逐策略删除
        & 负向秩关系保持；Holm 后 $p=0.100$，独立主效应未确认 \\

        SOC 区间
        & 替代切分与探索模型
        & 中高 SOC 方向反复为正但均未显著 \\

        \bottomrule
    \end{tabularx}
\end{table}
这些扰动不改变表\ref{tab:p2-conclusion-first}的证据等级：总体差异由S8驱动，$Q_1$仅为方向稳定的优先信号，SOC结论仍属于弱证据。

\subsection{问题三误差分析}
短期结论采用2轮$\times$7折策略约束Round-robin Outer验证，共14个外层折，40块训练电池均恰好进入验证2次。嵌套选择式Hybrid的折级平均RMSE、MAE分别为0.0480和0.0278个百分点，较Inner选择的最佳单模型改善5.7\%；按全部验证轨迹加权的PICP为98.8\%，TICP为92.5\%，MPIW为0.531个百分点。PICP评价单点覆盖，TICP要求50点轨迹全部覆盖，二者不可替代；较旧方案更宽的经验带体现了对整轨迹不确定性的保守处理。表\ref{tab:p3length}固定PCA--Ridge架构，只检验输入长度，不能与Hybrid主误差直接合并。

EOL没有真实阈值标签，不采用短期误差解释。联合抽取同一训练电池的完整参数校准向量可保留参数相关结构，三类模型的2.5\%--97.5\%包络则描述参数与模型形式不确定性，但不具备95\%覆盖保证。外推指数越大，中心值越应谨慎解释。

\subsection{问题四决策稳健性}\label{sec:p4-sensitivity}
最终Guard的策略级Outer LOPO中，时间与健康RMSE分别为0.4379 min和0.9060 log-rate，$R^2$分别为$-0.419$和$-0.432$。负$R^2$仍表明连续代理的整策略泛化不足，因此外层RMSE继续作为不确定性下限，连续解仅用于候选生成。

可信阈值取$0.75\delta_0$--$1.25\delta_0$、风险尺度取四种情景时，Balanced候选保持不变；1200次Monte Carlo中Fast、Balanced和Life区域出现非支配点的概率分别为58.4\%、95.7\%和90.2\%。稳定性主要来自共同的Outer-LOPO误差下限，而非代理已达到高精度。实测策略的20,000次电池级Bootstrap进一步表明速度端S4/S5/S7存在前沿归属不确定性，故以S5作为代表而非宣称唯一最优；S9在健康端的Pareto稳定率为83.7\%。


\section{模型评价与推广}
\subsection{模型优点}
\begin{itemize}
  \item \textbf{口径正确：}首先识别全部寿命事件右删失，避免把观测长度误作真实寿命，这是后续结论可信的前提。
  \item \textbf{证据分层：}正式推断、方向性探索和机理辅助具有明确边界，避免用弱模型放大结论。
  \item \textbf{防泄漏且可复现：}按电池分组验证，每问设置独立程序，中间表和日志均可追溯。
  \item \textbf{风险进入决策：}EOL报告跨模型包络；问题四以Outer LOPO误差约束连续候选，并用报告精度$\varepsilon$-Pareto及电池级Bootstrap区分“代表策略”与“唯一最优”。
\end{itemize}

\subsection{模型局限}
\begin{itemize}
  \item 只有9种策略，问题二同批次主分析仅有6个独立策略点；$C_1,Q_1,C_2$不是正交实验设计，无法识别严格因果主效应与交互效应。
  \item 数据只覆盖早期SOH，尚未观察拐点、加速衰减或80\%终点；轻微衰减电池的EOL只能看作风险提示。
  \item 温度是平均量，缺乏电芯内部峰值温度、负极电位与析锂直接观测，机理解释仍属于与文献一致的推断。
\end{itemize}

\subsection{改进与推广}
下一轮实验可围绕Balanced候选区域和S9健康端设置局部响应面，独立改变$Q_1$与$C_2$，增加重复电池并记录峰值温度、阻抗谱和析锂信号。获得完整寿命事件后，可采用加速失效时间或时变生存模型，并以安全约束贝叶斯优化安排后续试验\cite{attia2020}。迁移至其他材料体系时需重新标定SOC、温度与EOL边界。


\section*{AI工具使用声明}
本参赛队在竞赛过程中使用了AI工具，主要用于代码框架生成与调试、统计分析方案讨论、图表排版、论文结构检查及语言润色。AI未用于生成、补造或修改实验数据；所有程序均由参赛队在本地环境运行，核心数据口径、模型选择、公式、数值结果、图表和结论均经人工审查与复核。详细工具型号、使用环节、提示方式、采纳修改及核验情况见支撑材料《AI工具使用详情.pdf》。

\addcontentsline{toc}{section}{参考文献}
\begingroup
\small
\bibliographystyle{unsrt}
\bibliography{book}
\endgroup


\clearpage
\appendix
\section{支撑材料文件列表}
\begin{longtable}{p{0.34\textwidth}p{0.58\textwidth}}
\toprule
文件或目录 & 说明\\
\midrule
\endfirsthead
\toprule 文件或目录 & 说明\\ \midrule
\endhead
README.md，requirements.txt & PyCharm工程说明、运行顺序与Python依赖\\
code/common.py & 数据读取、raw/clean副本、右删失审计、稳健指标、统计检验及统一绘图函数\\
code/problem1.py & 问题一完整求解、表格和图形生成程序\\
\path{results/problem1/task1_battery_inventory.csv} & 49块电池的策略、寿命下界、充电时间与SOH指标明细\\
\path{results/problem1/baseline_normalization_sensitivity.csv} & 原始SOH及三个稳定基准窗的策略排序敏感性\\
\path{results/problem1/task4_charge_time_bootstrap.csv} & S8与S9两种充电时间差的电池级Bootstrap区间\\
code/problem2.py & 问题二五任务程序：精确检验、探索性Ridge/LOPO、SOC分区与机理一致性\\
\path{results/problem2/q2_soc_partition_sensitivity.csv} & SOC总水平、高低对比和两组替代切分结果\\
\path{results/problem2/q2_quasi_matched_pairs.csv} & S4/S5、S6/S7、S8/S9局部准匹配对比\\
code/q3\_solve\_base.py & 问题三统一的数据清洗、特征构造、候选模型与EOL基础函数\\
code/problem3\_v3.py & 问题三最终求解脚本：策略约束嵌套验证、短期Hybrid、TICP及EOL完整参数向量联合抽样\\
code/q3\_figures\_v3.py & 问题三图件的统一编号、多格式导出与质量检查程序\\
\path{results/problem3_v3/data/feature_group_permutation_summary_v3.csv} & 五组特征置换的Outer验证汇总结果\\
results/problem3\_v3 & 当前脚本输出的CSV、配置、运行记录、12幅问题三图件及图形质量检查记录\\
code/problem4.py & 问题四修订程序：Q1时间血缘、Guard Outer LOPO、可信域、风险Pareto与电池级Bootstrap\\
\path{results/problem4/results/q1_q4_data_lineage_check.csv} & S4--S9时间和退化率在Q1、Q4间逐项同源核验\\
\path{results/problem4/results/final_guard_outer_LOPO_summary.csv} & 时间/健康最终Guard的策略级Outer-LOPO误差\\
\path{results/problem4/results/continuous_representative_candidates.csv} & Fast、Balanced、Life连续代表候选及风险目标\\
\path{results/problem4/results/observed_policy_pareto_bootstrap.csv} & 六个实测策略在20,000次电池级Bootstrap中的经验Pareto频率\\
\path{results/problem4/results/observed_policy_pareto_epsilon_check.csv} & 未取整点Pareto与$\varepsilon_T=0.001$ min实践Pareto对照\\
\path{results/problem4/results/practical_recommendation.csv} & S5速度/均衡代表、S9健康优先与连续候选证据分级\\
code/overall\_flowchart.py & 论文整体技术路线图生成程序\\
code/qa\_checks.py & 核心数值、预测规模、推荐结论与图件一致性检查程序\\
results/problem1--problem4 & 全部中间表、检验结果、预测结果与模型文件\\
texfile/figures & 论文图件；问题四采用交付包的300 dpi PNG，其余主要图件同时保留SVG、PDF与TIFF\\
AI工具使用详情.pdf & 工具、环节、提示方式、采纳修改及人工核验说明\\
\bottomrule
\end{longtable}

\section{AI工具使用声明}
本参赛队在竞赛过程中使用了AI工具，主要用于代码框架生成与调试、统计分析方案讨论、图表排版、论文结构检查及语言润色。AI未用于生成、补造或修改实验数据；所有程序均由参赛队在本地环境运行，核心数据口径、模型选择、公式、数值结果、图表和结论均经人工审查与复核。工具型号、具体使用环节、主要提示方式、采纳修改及人工核验情况见支撑材料《AI工具使用详情.pdf》。

\section{程序复现说明}
完整Python源程序、逐问中间结果、CSV/JSON数值表及正文图件均随支撑材料压缩包交付。解压后按README.md配置环境：问题一、二、四运行各自脚本，问题三依次运行\path{code/problem3_v3.py}与\path{code/q3_figures_v3.py}，最后运行\path{code/qa_checks.py}，即可复核正文核心数值、预测规模、推荐结论与图件完整性。

竞赛提供的\path{battery_summary.csv}与\path{cycle_train.csv}不在支撑材料中重复收录。复现时应保持原文件名，将两文件放入工程\path{source}目录下的同一子目录。全部活动程序采用相对路径定位数据和输出目录；问题三的\path{q3_solve_base.py}必须与\path{problem3_v3.py}置于同一\path{code}目录。

\clearpage
\section{关键源程序代码（节选）}
本节仅列出足以说明模型实现、验证边界和不确定性处理的关键代码，不重复打印绘图程序和工程辅助代码。全部可运行源程序及其相对目录结构见支撑材料压缩包；历史版本\path{problem3.py}、\path{problem3_v2.py}及问题四旧版结果不参与当前论文复现。

\subsection{问题一：综合退化指标与充电时间差Bootstrap}
\lstinputlisting[style=Python,firstline=39,lastline=76,firstnumber=39,caption={综合退化指标权重敏感性}]{code/problem1.py}
\lstinputlisting[style=Python,firstline=177,lastline=201,firstnumber=177,caption={S8与S9充电时间差的电池级Bootstrap}]{code/problem1.py}

\subsection{问题二：精确推断与整策略留出验证}
\lstinputlisting[style=Python,firstline=168,lastline=200,firstnumber=168,caption={精确Mann--Whitney--Holm与精确Spearman检验}]{code/problem2.py}
\lstinputlisting[style=Python,firstline=331,lastline=388,firstnumber=331,caption={策略级LOPO Ridge评估}]{code/problem2.py}

\subsection{问题三：嵌套选择、EOL联合抽样与组置换}
\lstinputlisting[style=Python,firstline=135,lastline=180,firstnumber=135,caption={仅在Outer训练集内执行的Inner模型选择}]{code/problem3_v3.py}
\lstinputlisting[style=Python,firstline=330,lastline=383,firstnumber=330,caption={EOL完整参数向量的联合残差抽样}]{code/problem3_v3.py}
\lstinputlisting[style=Python,firstline=437,lastline=469,firstnumber=437,caption={五组特征的留出集联合置换重要性}]{code/problem3_v3.py}

\subsection{问题四：最终Guard与风险Pareto优化}
程序内部历史变量 \texttt{D}、\texttt{muD}、\texttt{sigD}、\texttt{F\_D}（及 \texttt{kD}、\texttt{sD}）对应正文健康响应 $Y$ 及其 Guard 目标（正文统一使用 $\mu_Y,\sigma_Y,F_Y$），不表示端点损失 $D_i$；为不影响真实代码复现，程序逻辑与变量名未作修改。
\lstinputlisting[style=Python,firstline=176,lastline=202,firstnumber=176,caption={时间与健康双分支Guard的Outer LOPO}]{code/problem4.py}
\lstinputlisting[style=Python,firstline=261,lastline=273,firstnumber=261,caption={可信域内风险修正Pareto与三类代表候选}]{code/problem4.py}




\end{document}
